% ============================ WORKSHEET 5 ==================================
\worksheetheading{5}{Atomic properties and periodicity}
\data{$h = \SI{6.62e-34}{\joule\second}$\quad $c = \SI{3.00e8}{\metre\per\second}$\quad
$1\ \mathrm{eV} = \SI{1.6e-19}{\joule}$\quad $N_A = \SI{6.022e23}{\per\mole}$\quad
$1\ \mathrm D = \SI{3.336e-30}{\coulomb\metre}$. Use Tables 5-1 to 5-4 of the handout (radii, ionisation
energies, bond energies, Pauling electronegativities). Calculator allowed.}

\wq{1}{Definitions}
Define: atomic radius, first ionisation energy, polarisability, electronegativity. For each, say
whether it is a property of an isolated atom or of an atom in a bond.
\begin{solution}
Atomic radius: distance from the nucleus to the outer edge of the electron cloud (radius of the highest
occupied AO); measured on atoms in molecules or solids.
First ionisation energy: minimum energy to remove the outermost electron from an atom in the gas phase,
$\mathrm X(\mathrm g)\to\mathrm X^+(\mathrm g)+\mathrm e^-$; isolated atom.
Polarisability: how easily the electron cloud of an atom is distorted by an electric field; isolated atom.
Electronegativity: ability of an atom \emph{in a bond} to attract the shared electrons.
\end{solution}

\wq{1}{Average atomic mass of chlorine}
Chlorine is 75.77\,\% $\iso{35}{Cl}$ (34.97~amu) and 24.23\,\% $\iso{37}{Cl}$ (36.97~amu). Compute its
average atomic mass and its molar mass.
\begin{solution}
$\bar m = 0.7577\times34.97 + 0.2423\times36.97 = 26.50 + 8.96 = 35.45$~amu; $M = 35.45\ \mathrm{g\,mol^{-1}}$.
\end{solution}

\wq{1}{Ranking radii}
Using Table 5-1, rank Na, Mg, K and Cl by increasing atomic radius and state the two trends that
explain the order.
\begin{solution}
Cl (99) $<$ Mg (160) $<$ Na (191) $<$ K (235~pm). The radius decreases across a period (Na $>$ Mg $>$ Cl,
period 3) and increases down a group (K $>$ Na, group 1).
\end{solution}

\wq{1}{Units of ionisation energy}
The ionisation energy of sodium is $496\ \mathrm{kJ\,mol^{-1}}$. Convert it to joules per atom and to
electronvolts per atom. Show that $1\ \mathrm{eV}$ per atom corresponds to $96.4\ \mathrm{kJ\,mol^{-1}}$.
\begin{solution}
$496\times10^3/6.022\times10^{23} = 8.24\times10^{-19}$~J per atom $= 8.24\times10^{-19}/1.6\times10^{-19}
= 5.15$~eV.\quad $1\ \mathrm{eV}\times N_A = 1.6\times10^{-19}\times6.022\times10^{23} = 9.64\times10^4\
\mathrm{J\,mol^{-1}} = 96.4\ \mathrm{kJ\,mol^{-1}}$ (96.5 with the exact value of $e$).
\end{solution}

\wq{2}{Mass versus atomic number}
(a) Compute the average atomic mass of potassium: $\iso{39}{K}$ 38.96~amu, 93.26\,\%;
$\iso{41}{K}$ 40.96~amu, 6.73\,\%. (b) Argon (39.95~amu) has $Z = 18$ and potassium $Z = 19$. What problem
would appear if elements were ordered by mass? How is it solved today?
\begin{solution}
(a) $0.9326\times38.96 + 0.0673\times40.96 = 36.33 + 2.76 = 39.09$~amu.\\
(b) By mass K (39.09) would come before Ar (39.95): K would fall under the noble gases and Ar under the
alkali metals, against their chemistry. The table is ordered by $Z$ (number of protons), which puts
Ar in group 18 and K in group 1.
\end{solution}

\wq{2}{Explaining the radius trends}
Explain, using the valence shell $n$ and the effective nuclear charge, why the radius decreases from
Li (157~pm) to F (64~pm) and increases from F to I (133~pm). Why can Ne (154~pm) not be compared with F?
\begin{solution}
Li to F: same valence shell ($n = 2$), but $Z_\mathrm{eff}$ increases because each added electron is in the
same shell and screens the added proton poorly; the cloud is pulled in. F to I: $n$ goes from 2 to 5,
the valence orbitals are much larger ($r\propto n^2$), while $Z_\mathrm{eff}$ barely changes. Ne's value
is a van der Waals radius (non-bonded atoms), always much larger than a covalent radius.
\end{solution}

\wq{2}{An isoelectronic series}
Rank $\ce{O^2-}$, $\ce{F-}$, $\ce{Na+}$ and $\ce{Mg^2+}$ by increasing radius. Justify. Compare also the
radius of $\ce{Na+}$ with that of Na, and of $\ce{Cl-}$ with that of Cl.
\begin{solution}
All have 10 electrons ($1\mathrm s^2 2\mathrm s^2 2\mathrm p^6$) but 8, 9, 11 and 12 protons. More protons for
the same electrons pull the cloud in: $\ce{Mg^2+} < \ce{Na+} < \ce{F-} < \ce{O^2-}$.
$\ce{Na+}$ is much smaller than Na (it loses its whole $n = 3$ shell); $\ce{Cl-}$ is larger than Cl (one more
electron, more repulsion, same nucleus).
\end{solution}

\wq{2}{The ionisation energies of period 2}
The values for period 2 are: Li 520, Be 900, B 801, C 1086, N 1402, O 1314, F 1681,
Ne 2081~$\mathrm{kJ\,mol^{-1}}$. (a) State the general trend and explain it. (b) Explain the two
exceptions with box diagrams.
\begin{solution}
(a) $E_I$ increases across the period: $Z_\mathrm{eff}$ increases and the radius decreases, so the
electron is held more strongly.\\
(b) Be $>$ B: B loses its single 2p electron (\obox{d}\;\obox{u,e,e}), higher in energy and screened by
2s, easier to remove than a 2s electron of Be (\obox{d}\;\obox{e,e,e}).
N $>$ O: N has a half-filled 2p (\obox{u,u,u}); O has one pair (\obox{d,u,u}), and the repulsion
between the two paired electrons makes one of them easier to remove.
\end{solution}

\wq{3}{Ionising light}
(a) Compute the longest wavelength of a photon able to ionise a sodium atom (496~$\mathrm{kJ\,mol^{-1}}$),
and the one for helium (2372~$\mathrm{kJ\,mol^{-1}}$). In which region of the spectrum are they?
(b) In a PES experiment, UV photons of 21.2~eV hit sodium atoms. What is the kinetic energy of the
ejected valence electrons?
\begin{solution}
(a) Na: $E = 8.24\times10^{-19}$~J (Q4), $\lambda = hc/E = 6.62\times10^{-34}\times3.00\times10^8/8.24\times10^{-19}
= 2.41\times10^{-7}$~m $= 241$~nm (UV).
He: $E = 2372\times10^3/6.022\times10^{23} = 3.94\times10^{-18}$~J, $\lambda = 50.4$~nm (far UV).\\
(b) $E_k = h\nu - E_I = 21.2 - 5.15 = 16.0$~eV.
\end{solution}

\wq{3}{Successive ionisation energies}
An element of period 3 has successive ionisation energies 738, 1451, 7733, 10\,543~$\mathrm{kJ\,mol^{-1}}$.
(a) Identify it and justify. (b) Why does each ionisation energy exceed the previous one? (c) Which ion
does this element form in its compounds?
\begin{solution}
(a) The huge jump comes between the 2nd and the 3rd ($1451\to7733$): 2 electrons are easy to remove,
the 3rd is a core electron (2p). Two valence electrons in period 3: magnesium, $[\mathrm{Ne}]\,3\mathrm s^2$.\\
(b) Each electron is removed from an increasingly positive ion: fewer electrons for the same nuclear
charge, less repulsion, the remaining electrons are held more tightly.\\
(c) $\ce{Mg^2+}$ ($[\mathrm{Ne}]$).
\end{solution}

\wq{3}{Abundance from the average mass}
Boron has two isotopes, $\iso{10}{B}$ (10.013~amu) and $\iso{11}{B}$ (11.009~amu), and an average atomic
mass of 10.81~amu. Compute the natural abundance of each isotope.
\begin{solution}
With $x$ the fraction of $\iso{10}{B}$: $10.013\,x + 11.009\,(1-x) = 10.81$, so
$x = (11.009 - 10.81)/(11.009 - 10.013) = 0.199/0.996 = 0.200$: 20.0\,\% $\iso{10}{B}$ and
80.0\,\% $\iso{11}{B}$.
\end{solution}

\wq{3}{How polar is H--Cl?}
The dipole moment of HCl is 1.08~D and its bond length 127~pm. (a) Compute the partial charge $\delta$
(as a fraction of $e$) carried by each atom. Which atom is $\delta-$? (b) Using Table 5-4, rank
H--F, H--Cl, H--Br and H--I by increasing polarity.
\begin{solution}
(a) $\delta e = \mu/d = 1.08\times3.336\times10^{-30}/127\times10^{-12} = 2.84\times10^{-20}$~C, so
$\delta = 2.84\times10^{-20}/1.6\times10^{-19} = 0.18$. Cl is more electronegative (3.16 vs 2.2): Cl is
$\delta-$, H is $\delta+$. The bond is about 18\,\% ionic.\\
(b) $\Delta\chi$: H--I $2.66 - 2.2 = 0.46$ $<$ H--Br 0.76 $<$ H--Cl 0.96 $<$ H--F 1.78.
\end{solution}

\wq{4}{Pauling electronegativity from bond energies}
Use Table 5-3 (H--H 436, F--F 158, H--F 567, I--I 151, H--I 299~$\mathrm{kJ\,mol^{-1}}$), the formula
$|\chi_A - \chi_B| = 0.102\sqrt{\Delta}$ with $\Delta = D_{AB} - \sqrt{D_{AA}D_{BB}}$, and $\chi_\mathrm H = 2.20$.
(a) Compute $\chi_\mathrm F$. (b) Compute $\chi_\mathrm I$ and compare with Table 5-4 (2.66). (c) Why is
the H--F bond much stronger than the average of H--H and F--F?
\begin{solution}
(a) $\sqrt{436\times158} = 262.5$; $\Delta = 567 - 262.5 = 304.5$; $0.102\sqrt{304.5} = 1.78$;
F is more electronegative than H: $\chi_\mathrm F = 2.20 + 1.78 = 3.98$ (matches Table 5-4).\\
(b) $\sqrt{436\times151} = 256.6$; $\Delta = 299 - 256.6 = 42.4$; $0.102\sqrt{42.4} = 0.66$;
$\chi_\mathrm I = 2.20 + 0.66 = 2.86$, 0.2 above the table. Pauling's values are fitted over many
molecules; one pair only gives an estimate, and small $\Delta$ values are very sensitive to errors.\\
(c) In the polar H--F bond the partial charges $\delta+$ and $\delta-$ attract each other: this ionic
contribution adds to the covalent bond energy. The more different the electronegativities, the larger
this extra energy $\Delta$.
\end{solution}

\wq{4}{The Allred--Rochow scale}
(a) Using Slater's rules (Unit 5, ``Going further''), show that $Z_\mathrm{eff} = 6.1$ for a valence
electron of Cl. (b) With $\chi = 0.359\,Z_\mathrm{eff}/r^2 + 0.744$ ($r$ in \AA) and the covalent radius
of Cl (99~pm), compute $\chi_\mathrm{AR}$(Cl) and compare with Pauling's 3.16. (c) Without calculation,
explain why F has a higher electronegativity than Cl on this scale.
\begin{solution}
(a) Cl $= 1\mathrm s^2\,2\mathrm s^2 2\mathrm p^6\,3\mathrm s^2 3\mathrm p^5$. For one 3p electron:
6 other $n = 3$ electrons $\times0.35 = 2.10$; 8 electrons of $n = 2$ $\times0.85 = 6.80$; 2 electrons of
$n = 1$ $\times1.00 = 2.00$; $\sigma = 10.9$ and $Z_\mathrm{eff} = 17 - 10.9 = 6.1$.\\
(b) $r = 0.99$~\AA: $\chi = 0.359\times6.1/0.99^2 + 0.744 = 2.23 + 0.74 = 2.98$. The two scales agree
within about 0.2 (they are built on different quantities: a force for Allred--Rochow, bond energies
for Pauling).\\
(c) F has a similar $Z_\mathrm{eff}$ (5.2 by Slater's rules) but a much smaller radius (64~pm): its
valence shell is $n = 2$. $Z_\mathrm{eff}/r^2$ is larger, so the bonding electrons are attracted
more strongly.
\end{solution}

\wq{4}{Synthesis: an unknown element}
Element X is in period 3. Its successive ionisation energies are 578, 1817, 2745 and
11\,577~$\mathrm{kJ\,mol^{-1}}$. (a) Identify X and write its configuration with the box diagram of its
valence shell. (b) Explain why the first ionisation energy of X is lower than that of Mg (738), although
X has one more proton. (c) Compare the radius, the electronegativity and the polarisability of X with
those of Mg and Si, using the tables. (d) Which ion does X form? Is it isoelectronic with a noble gas?
\begin{solution}
(a) The jump is after the 3rd electron (2745 $\to$ 11\,577): 3 valence electrons in period 3: aluminium,
$[\mathrm{Ne}]\,3\mathrm s^2 3\mathrm p^1$, \obox{d}\;\obox{u,e,e}.\\
(b) Al loses its single 3p electron, which is higher in energy and partly screened by the 3s pair;
Mg loses a 3s electron, more penetrating and more tightly held (same anomaly as Be/B).\\
(c) Radius: Mg 160 $>$ Al 143 $>$ Si 118~pm (decreases across the period). Electronegativity:
Mg 1.31 $<$ Al 1.61 $<$ Si 1.90 (increases). Polarisability follows the radius: Mg $>$ Al $>$ Si.\\
(d) $\ce{Al^3+}$, $[\mathrm{Ne}]$: isoelectronic with neon.
\end{solution}
